% Speedup of the n:1 over the 1:1 translation of VIP, per host and vector width, drawn from
% the tables that plot.py writes to data/ (median, min and max of the per-run ratios).
\documentclass[9pt]{amsart}
\usepackage{pivotfig}

\pgfplotstableread[col sep=comma]{data/sapphirerapids.csv}\vipSpr
\pgfplotstableread[col sep=comma]{data/znver4.csv}\vipZen
\pgfplotstableread[col sep=comma]{data/a64fx.csv}\vipAfx

% Every kernel group gets the same slot width, and one doubling the same height.
\pgfmathsetlengthmacro{\vipSlot}{20pt}
\pgfmathsetlengthmacro{\vipOctave}{18pt}
\pgfmathsetlengthmacro{\vipBar}{0.27*\vipSlot}
\pgfmathsetlengthmacro{\vipShift}{0.29*\vipSlot}
\pgfplotsset{
  vip panel/.style={
    log origin=0, clip=false, scale only axis, x=\vipSlot, enlarge x limits={abs=0.6},
    ymin=0.8, ymax=6, height={ln(6/0.8)/ln(2)*\vipOctave},
    axis lines*=left, axis line style={gray1, line width=0.4pt},
    xtick=data, xtick align=outside, xticklabels from table={#1}{kernel},
    xticklabel style={rotate=45, anchor=north east, inner sep=1pt, font=\footnotesize, text=gray1},
    ytick={1,2,4,8}, yticklabels={},
    tick style={gray1, line width=0.4pt}, ymajorgrids,
    major grid style={gray1!25, line width=0.2pt},
    yticklabel style={font=\footnotesize, text=gray1},
    ybar, bar width=\vipBar,
  },
  vip left/.style={yticklabel={\pgfmathparse{exp(\tick)}\pgfmathprintnumber[fixed]{\pgfmathresult}},
    ylabel={speedup}, ylabel style={font=\footnotesize, text=gray1}},
}
\tikzset{
  vip 128/.style={draw=magenta1, fill=magenta2, line width=0.7pt},
  vip 256/.style={draw=teal1, fill=teal2, line width=0.7pt},
  vip 512/.style={draw=blue1, fill=blue2, line width=0.7pt},
}
% Draws the bars of one width, at the median over runs.
\newcommand{\vipBars}[3]{%
  \addplot[vip #2, bar shift=#3] table[x expr=\coordindex, y=med#2] {#1};
}
% Draws one panel: the three widths, the line at 1, and the title.
\newcommand{\vipPanel}[2]{%
  \vipBars{#1}{128}{-\vipShift}
  \vipBars{#1}{256}{0pt}
  \vipBars{#1}{512}{\vipShift}
  \draw[gray1, line width=0.4pt] ({rel axis cs:0,0}|-{axis cs:0,1}) -- ({rel axis cs:1,0}|-{axis cs:0,1});
  \node[font=\footnotesize, text=gray1, anchor=base] at ([yshift=5pt]rel axis cs:0.5,1) {#2};
}

\begin{document}
\begin{tikzpicture}
\begin{axis}[ymode=log, vip panel=\vipSpr, vip left, name=spr]
  \vipPanel{\vipSpr}{Intel Sapphire Rapids}
\end{axis}
\begin{axis}[ymode=log, vip panel=\vipZen, name=zen,
    at={($(spr.north east)+(6pt,0)$)}, anchor=north west]
  \vipPanel{\vipZen}{AMD Zen 4}
\end{axis}
\begin{axis}[ymode=log, vip panel=\vipAfx, name=afx,
    at={($(zen.north east)+(6pt,0)$)}, anchor=north west]
  \vipPanel{\vipAfx}{Fujitsu A64FX}
\end{axis}
\coordinate (legend) at ($(spr.south west)!0.5!(afx.south east)$);
\matrix[anchor=north, yshift=-6pt, column sep=2pt,
    nodes={font=\footnotesize, text=gray1, inner sep=0pt, anchor=base west}]
    at (legend |- current bounding box.south) {
  \draw[vip 128] (0pt,-0.6pt) rectangle (5pt,4.4pt); & \node {SSE/SVE (128 bit)}; &[5pt]
  \draw[vip 256] (0pt,-0.6pt) rectangle (5pt,4.4pt); & \node {AVX2/SVE (256 bit)}; &[5pt]
  \draw[vip 512] (0pt,-0.6pt) rectangle (5pt,4.4pt); & \node {AVX-512/SVE (512 bit)}; \\
};
\end{tikzpicture}
\end{document}
